\newpage
\section*{References}
\begingroup
\setlength{\parindent}{-14pt}\setlength{\leftskip}{14pt}\setlength{\parskip}{3.2pt}\small
Sources marked (abstract) were seen only as an abstract. Some figures came through an automated reader and still need a human check; \texttt{report/literature/} lists them.

Anthropic. (n.d.). \emph{Use Claude's chat search and memory to build on previous context}. Claude Help Center. Retrieved November 5, 2026, from \url{https://support.claude.com/en/articles/11817273}

Bastani, H., et al. (2025). Generative AI without guardrails can harm learning. \emph{PNAS, 122}(26), e2422633122. \url{https://doi.org/10.1073/pnas.2422633122}

Bensal, S., Magnuson, A., Balagopalan, A., \& Bikel, D. M. (2026). \emph{Recalling too well: Sycophancy evaluation and mitigation in memory-augmented models} (Version 2). arXiv:2606.10949. All authors are employed by WRITER; the preprint states no competing interests.

Bergson, Z., Vassall, S. G., Wright, A., McCoy, A. B., Schafer, K. M., Achee, M. C., \& Sheffield, J. M. (2026). Characterizing artificial intelligence psychosis in a large medical setting. \emph{JAMA Psychiatry}. \url{https://doi.org/10.1001/jamapsychiatry.2026.2923}

Blank, C., Ying, Z., Potts, C., Hase, P., \& Huang, J. (2026). \emph{Sycophantic agreement transfers with neutral data via contrastive preference optimization}. arXiv:2608.31079.

Bo, J. Y., Kazemitabaar, M., Deng, M., Inzlicht, M., \& Anderson, A. (2026). Invisible saboteurs: Sycophantic LLMs mislead novices in problem-solving tasks. CHI 2026. arXiv:2510.03667.

Butterfield, B., \& Metcalfe, J. (2001). Errors committed with high confidence are hypercorrected. \emph{Journal of Experimental Psychology: Learning, Memory, and Cognition, 27}(6), 1491--1494 (abstract).

Campus Technology. (2025, July 30). \emph{New ChatGPT study mode guides students through questions}. \url{https://campustechnology.com/Articles/2025/07/30/New-ChatGPT-Study-Mode-Guides-Students-Through-Questions.aspx}

Chen, S., Ying, Z. J., Smith, L. R., Wertheimer, J., \& Shapira, N. (2026). \emph{Tracing mechanisms of sycophantic agreement in language models}. arXiv:2609.35822.

De Marez, V., De Bruyne, L., \& Daelemans, W. (2026). \emph{Decomposing factual sycophancy in language models: How size and instruction tuning shape robustness}. arXiv:2606.06306.

Du, K., K\"umpel, C., Wastl, M., \& Warstadt, A. (2026). It's not what you say, it's how you say it: Evaluating LLM responses to expressions of belief. In \emph{Proceedings of ACL 2026}. \url{https://aclanthology.org/2026.acl-long.142.pdf}

Elish, M. C. (2019). Moral crumple zones: Cautionary tales in human-robot interaction. \emph{Engaging Science, Technology, and Society, 5}, 40--60. \url{https://doi.org/10.17351/ests2019.260}

Federal Rules of Evidence, Rule 803(3) (United States). Legal Information Institute. Retrieved November 5, 2026, from \url{https://www.law.cornell.edu/rules/fre/rule_803}

Fricker, M. (2007). \emph{Epistemic injustice: Power and the ethics of knowing}. Oxford University Press. Unread by the authors; cited as quoted in Lackey (2018).

Fricker, M. (2017). Evolving concepts of epistemic injustice. In I. J. Kidd, J. Medina, \& G. Pohlhaus (Eds.), \emph{The Routledge handbook of epistemic injustice} (pp. 53--60). Routledge.

Genadi, R., Nwadike, M., Mukhituly, N., Alquabeh, H., Hiraoka, T., \& Inui, K. (2026). \emph{Sycophancy hides linearly in the attention heads}. arXiv:2601.16644.

Google. (n.d.). \emph{Get personalization with memory of your past Gemini chats}. Gemini Apps Help. Retrieved November 7, 2026, from \url{https://support.google.com/gemini/answer/16598469}

Hannoon, H., Zhao, A., Narayan, M., Goyal, S., \& Sheth, I. (2026). \emph{Mitigating over-personalization in LLMs via structured memory}. arXiv:2608.08300.

Hendrycks, D., Burns, C., Basart, S., Zou, A., Mazeika, M., Song, D., \& Steinhardt, J. (2020). \emph{Measuring massive multitask language understanding}. arXiv:2009.03300.

Hong, J., Byun, G., Kim, S., Shu, K., \& Choi, J. D. (2025). \emph{Measuring sycophancy of language models in multi-turn dialogues}. arXiv:2505.23840.

Jain, S., Park, C., Viana, M., Wilson, A., \& Calacci, D. (2026). Interaction context often increases sycophancy in LLMs. In \emph{Proceedings of the 2026 CHI Conference on Human Factors in Computing Systems}. \url{https://doi.org/10.1145/3772318.3791915}

Kelly, T. (2005). The epistemic significance of disagreement. In T. S. Gendler \& J. Hawthorne (Eds.), \emph{Oxford studies in epistemology} (Vol. 1). Oxford University Press.

Kim, T. M., Luo, L., Kim, S. E., Manrai, A. K., Topol, E., \& Rajpurkar, P. (2026). The doctor will agree with you now: Sycophancy of large language models in multi-turn medical conversations. HeaLing workshop. \url{https://aclanthology.org/2026.healing-1.2.pdf}

Lackey, J. (2018). Credibility and the distribution of epistemic goods. In K. McCain (Ed.), \emph{Believing in accordance with the evidence} (pp. 145--168). Springer.

LearnLM Team \& Eedi. (2025). \emph{AI tutoring can safely and effectively support students: An exploratory RCT in UK classrooms}. arXiv:2512.23633. Authors include Google staff.

Luz de Araujo, P. H., \& Roth, B. (2024). \emph{Helpful assistant or fruitful facilitator? Investigating how personas affect language model behavior}. arXiv:2407.02099.

Ma, H., Zou, H. P., Li, C., Ma, E., Su, Y., \& Yu, P. S. (2026). \emph{Sycophancy suppression can impair rational updating}. arXiv:2608.26511.

Mao, X., Zhao, L., Wang, L., Qian, R., Huang, Q., Wang, W., Han, B., Zheng, X., \& Wang, C. (2026). \emph{Agents don't just agree, they remember: Benchmarking persistent sycophancy in self-improving personal agents}. arXiv:2607.10526.

Ministry of Education, Singapore. (2026). \emph{AI usage in schools} (parliamentary reply, May 6), and \emph{Breakdown of Primary 3 to 6 students receiving private tuition} (parliamentary reply, February 12). \url{https://www.moe.gov.sg/news/parliamentary-replies}

Morrin, H., et al. (2025). \emph{Delusions by design?} PsyArXiv cmy7n, version 5. \url{https://osf.io/preprints/psyarxiv/cmy7n_v5}

Nan, X., Wang, Y., \& Thier, K. (2022). Why do people believe health misinformation and who is at risk? \emph{Social Science \& Medicine, 314}, 115398 (abstract).

National Institute of Education. (2024, July 29). \emph{MOE's newest AI tools and how schools are using them}. \url{https://ntu.edu.sg/nie/news-events/news/detail/moe-s-newest-ai-tools-and-how-schools-are-using-them}

Nguyen, C. T. (2020). Echo chambers and epistemic bubbles. \emph{Episteme, 17}(2), 141--161. \url{https://doi.org/10.1017/epi.2018.32}

NST2046. (2026, October 20). \emph{Who learns? What happens to learning when the answer arrives before the mistake} [Lecture slides]. National University of Singapore.

Pan, D. S., Huang, J. H., Lee, M. H., et al. (2016). Knowledge, attitudes and practices towards antibiotic use in upper respiratory tract infections among patients seeking primary health care in Singapore. \emph{BMC Family Practice, 17}(1), 148.

Pew Research Center. (2019). \emph{What Americans know about science}.

Pew Research Center. (2025). \emph{34\% of U.S. adults have used ChatGPT, about double the share in 2023}.

Ping, K., \c{C}ar{\i}k, B., Wohn, C., Ding, X., Wang, T., \& Rho, E. (2026). \emph{Why LLMs give in: Conversational factors and reasoning behind medical sycophancy}. arXiv:2608.01017.

Pugh, S. L., Yang, E., Sutherland, A. M., \& Breschi, A. (2026). \emph{Detecting clinical discrepancies in health coaching agents} (Version 1). arXiv:2604.27045. All authors are employed by Verily Health.

Pulipaka, S., Chen, O., Sharma, M., Bajwa, T. S., Raina, V., \& Sheth, I. (2026). \emph{PersistBench: When should long-term memories be forgotten by LLMs?} arXiv:2602.01146.

Sadler, P. M., Sonnert, G., Coyle, H. P., Cook-Smith, N., \& Miller, J. L. (2013). The influence of teachers' knowledge on student learning in middle school physical science classrooms. \emph{American Educational Research Journal, 50}(5), 1020--1049.

Schroeder, N. L., \& Kucera, A. C. (2022). Refutation text facilitates learning. \emph{Educational Psychology Review, 34}(2), 957--987.

Shapira, I., Cheng, M., \& Orgad, H. (2026). \emph{Is sycophancy a form of in-context updating?} arXiv:2610.28325.

Sharma, M., et al. (2023). \emph{Towards understanding sycophancy in language models}. arXiv:2310.13548.

Singapore Evidence Act 1893, sections 32 and 32C. Version in force on November 5, 2026. \url{https://sso.agc.gov.sg/Act/EA1893}

Soon, C., \& Goh, S. (2021). \emph{Singaporeans' susceptibility to false information} (IPS Exchange Series No. 19). Institute of Policy Studies.

Straits Times, via Asia News Network. (2024). \emph{Spending by Singapore families on private tuition rises to \$1.8 billion in 2023}. \url{https://asianews.network/?p=185968}

Sultan, M., Tump, A. N., Ehmann, N., et al. (2024). Susceptibility to online misinformation: A systematic meta-analysis of demographic and psychological factors. \emph{PNAS}. \url{https://doi.org/10.1073/pnas.2409329121}

Suzgun, M., Gur, T., Bianchi, F., Ho, D. E., Icard, T., Jurafsky, D., \& Zou, J. (2024). \emph{Belief in the machine: Investigating epistemological blind spots of language models}. arXiv:2410.21195.

Willison, S. (2025, May 2). \emph{Expanding on what we missed with sycophancy}. \url{https://simonwillison.net/2025/May/2/what-we-missed-with-sycophancy/}. Reproduces OpenAI's post.

Wu, K., Wu, E., \& Zou, J. (2024). \emph{ClashEval: Quantifying the tug-of-war between an LLM's internal prior and external evidence}. arXiv:2404.10198.

Xiang, Z., Chen, Z., Tang, Y., Wei, Z., Ning, R., Lin, Y., Zhang, Q., \& Su, J. (2026). \emph{MemSyco-Bench: Benchmarking sycophancy in agent memory} (Version 3). arXiv:2607.01071.

Yao, J., Parihar, S., Bui, N. T., Savova, G. K., \& Wang, L. (2026). \emph{Locating sycophancy across training stages in open language models}. arXiv:2610.37788.

Zhan, Q., Zhang, R., Guo, S., Zhao, L., \& Liu, Z. (2026). \emph{When memory becomes authority: Benchmarking authority collapse at the memory consolidation boundary}. arXiv:2608.01679.
\par\endgroup
